\scanpage{167}
\uline{Compatibilité combinée de $\varepsilon,\vee,\Delta$}, On veut que
pour $V\in K^{\underline U}$, l'hom naturel $V'\otimes V\longrightarrow k$
soit un $\underline U$-hom, quand $\underline k$ est considéré comme
$U$-module trivial (si on n'a pas supposé $\Delta$
symétrique, il faut supposer la $=$ chose pour
$V\otimes V'\longrightarrow k$). Ceci s'exprime par
\[ \underline{P(\vee\otimes1)\Delta=1}\ \text{et}\ \underline{P(\deleted{\unsure{(}}1\otimes\vee)\Delta=1} \]
(ces deux sont équivalentes si $\Delta$ est symétrique)

Si \uline{tous} ces axiomes sont vérifiés, alors
tout hom. fonctoriel classique construit sur
des $\underline U$-modules permis est permis.

\deleted{On peut \illegible{3} \unsure{qui soit} \unsure{tout}}

\noindent\rule{2.5cm}{0.4pt}

Soient $\underline U,\underline V$ deux ``$K$-algèbres'' avec
hom diagonal $\Delta$, $\varphi$ un hom \inserted{\uline{continu}} $\underline V\longrightarrow\underline U$
i.e. $\overline\varphi$ un foncteur permis $K^{\underline V}\longrightarrow K^{\underline U}$.
\uline{Pour que $\overline\varphi$ soit compatible avec les}
\uline{produits tensoriels} (i.e. si $V,W\in K^{\underline U}$,
alors sur $V\otimes W$ la structure de $\underline V$-module
associée est le produit \unsure{des st|deux st} à $\underline\varphi(\unsure{W|V})$ et $\underline\varphi(\unsure{w|W})$) il faut et il suffit
que $\varphi$ soit compatible avec $\Delta$.

% archive sheet number: 166
